\subsection{Computing the finite lower bound}\label{sec:recipe}
To convert Theorem~\ref{thm:finite-folded} into an explicit finite diagonal lower bound, we allocate the \(N\) terms across \(r\) fold stages whose main contributions share the same quadratic order, summing to the extra logarithmic factor. The allocation assigns \(p\) odd powers of \(y\) to \(P\), and \(X_*\) odd powers of \(x\) and \(d\) positive even powers of \(y\) to \(Q\). All logarithms in the allocation have base two. There are \(r\) fold stages. At stage \(k\), each of \(Y_k\) rows uses clusters of \(q_{r-k}\) wells, up to the common size limit \(q_{\max}\). For \(N\ge48\) set
\begin{align}
 r&=\max\{1,\lfloor\log_2 N\rfloor-7\},\quad R=2^r,\quad
 q_{\max}=2^{\lceil\log_2\max(2,r)\rceil},\notag\\
 X_*&=2q_{\max}\left\lfloor\frac{\min(\lfloor N/2\rfloor,N-18R)}{2q_{\max}}\right\rfloor,\qquad
 s=X_*/q_{\max},\quad T=N-X_*,\notag\\
 p&=\max(12R,\lfloor T/2\rfloor),\qquad d=T-p. \label{eq:allocation}
\end{align}
Use lengths \(\ell_i=i+3\), endpoints \(b_0=0,\ b_i=i(i+7)/2\), and \(I=\max\{i:b_i\le r\}\). At stage \(0\le k<r\) write \(j=r-k\) and compute
\begin{align}
 e_j&=2^{\max(1,\#\{i\ge1:b_i\le j\})},&
 Y_k&=2^{k+1}\prod_{\substack{1\le i\le I\\r-b_{i-1}\le k}}(1-4/2^{\ell_i}),\notag\\
 q_j&=\min\!\left(q_{\max},2^{\max(1,\lfloor(\log_2e_j)/2\rfloor)}\right),&
 \alpha_k&=2e_j,\notag\\
 M_k&=\left\lfloor\frac p{2^{k+1}}\right\rfloor+
       \left\lfloor\frac{d+2^k}{2^{k+1}}\right\rfloor,&
 L_k&=\max(0,M_k-\alpha_k). \label{eq:recipe-stages}
\end{align}
With the per-cluster count \(B_{q,\alpha}\) from \eqref{eq:Bqalpha}, the allocation gives
\begin{equation}\label{eq:finite-main}
 \mathcal F(N)=\sum_{k=0}^{r-1}Y_k\,\frac{X_*}{q_{r-k}}\,
          \max_{0\le\ell\le L_k}B_{q_{r-k},\alpha_k}(\ell).
\end{equation}
At each stage the maximum selects how many of the available modes to use. In a tie choose the largest maximizing prefix. For this allocation \(\alpha_k\ge2q_j\). Advancing \(\ell\ge1\) by \(\alpha_k\) increases \(b,c\) by one and two, respectively, and cannot reduce any of the three positive parts. Thus only the last \(\alpha_k\) prefixes and the zero case need to be considered.
The last stage \(k=r-1\) has \(q_1=2\) and \(\alpha_k=4\). Since \(T\ge18R\), its two floors give \(M_k>T/R-3/2\ge 33/2\); hence \(L_k\ge13\). The prefix \(\ell=8\) gives \(B_{2,4}(8)=2\) and the last-stage term in \eqref{eq:finite-main} is positive. This proves \(\mathcal F(N)>0\) for every \(N\ge48\).

For comparison at the same exact \(N\) and cap \(N-1\), write
\[
 \mathcal E(N)=\lfloor(N-2)^2/4\rfloor,\qquad
 \mathcal S(N)=(N-4)^2-\mathbf{1}_{\,2\mid N}.
\]
These are the elementary and stationary-turn counts of Propositions~\ref{thm:allocation} and~\ref{prop:stationary-budget}. For separate strips (Proposition~\ref{prop:strip-product}), put \(k=\lfloor N/2\rfloor\). For every feasible depth \(i\ge1\) take the largest integer \(b\ge4\) with \(s=2^i b\) satisfying
\[
 k+2^{i+1}+s-1\le N,\qquad
 \max(2k-1,2^{i+2}-1,2s-2)\le N-1 .
\]
The separate-strip count \(\mathcal T(N)\) is the maximum of \(2kc_i\) over these depths, where \(c_i=is/2-2^i+1\). The maximum here fixes the number \(k\) of horizontal strips. The degree of the folded field is \(D_{\rm f}=\max(2p-1,2X_*-1,2d)\).

Table~\ref{tab:fold} compares the finite formulas. The exact integer scan in \cite{LuReplay} over \(48\le N\le65{,}536\), with \(\mathcal S\) taken in its odd or even form according to the parity of \(N\), gives strict inequalities and no ties:
\begin{align*}
 \mathcal F(N)&<\mathcal E(N) &&(48\le N\le511),\\
 \mathcal E(N)<\mathcal F(N)&<\mathcal S(N) &&(512\le N\le8191),\\
 \mathcal F(N)&>\mathcal S(N) &&(8192\le N\le65{,}536).
\end{align*}
Both changes occur where the number of folds \(r=\max(1,\lfloor\log_2N\rfloor-7)\) increases, from \(1\) to \(2\) at \(N=512\) and from \(5\) to \(6\) at \(N=8192\). For the same two formulas \(\mathcal F(N)>N^2>\mathcal S(N)\) holds analytically for every \(N\ge 2^{15}=32{,}768\). The analytic range overlaps the finite scan. Indeed, for \(N\ge 2^{15}\) one has \(r\ge 8\), \(N\ge 128R\), and hence \(N-18R>\lfloor N/2\rfloor\). On each of the last eight stages \(k=r-8,\dots,r-1\) (\(j=r-k\in\{1,\dots,8\}\)), \(e_j=q_j=2\) and \(\alpha_k=4\). By \eqref{eq:ej-Yk} and \(1-\prod(1-\delta_i)\le\sum\delta_i\), every stage satisfies
\[
 \frac{Y_k}{2^{k+1}}
   = \prod_{\substack{2\le i\le I\\ r-b_{i-1}\le k}}\left(1-\frac{1}{2^{i+1}}\right)
   \ge 1-\sum_{i=2}^\infty\frac{1}{2^{i+1}} = \frac{3}{4}.
\]
For \(q=2,\ \alpha=4\), \eqref{eq:Bqalpha} gives \(b\le(L+1)/4\) and \(B_{2,4}(L)\ge 2J+2(J-b)=4L-10b-12\ge\tfrac{3}{2}L-\tfrac{29}{2}\) for \(L\ge1\). The final affine bound also holds at \(L=0\), where \(B_{2,4}(0)=0\). Since \(M_k>T/2^{k+1}-3/2\), we have \(L_k\ge M_k-4>T/2^{k+1}-11/2\), so \(B_{2,4}(L_k)>\tfrac{3}{2}T/2^{k+1}-\tfrac{91}{4}\). Multiplying by \(Y_k X_*/2=(Y_k/2^{k+1})\,2^k X_*\) and summing over \(k=r-8,\dots,r-1\) (using \(Y_k/2^{k+1}\in[\tfrac{3}{4},1]\) and \(\sum_{k=r-8}^{r-1}2^k<R\le N/128\le T/64\)) yields
\[
 \mathcal F(N)
   > \frac{9}{2}X_*T - \frac{91}{4}X_*R
   \ge \left(\frac{9}{2}-\frac{91}{256}\right)X_*T
   = \frac{1061}{1024}\,(4X_*T).
\]
Because \(q_{\max}<2r\) gives \(0\le N-2X_*\le 4q_{\max}-1<8r\le r\,2^{-r}N/16\le N/512\), we have \(4X_*T=N^2-(N-2X_*)^2>(1-2^{-18})N^2\), whence \(\mathcal F(N)>\frac{1061}{1024}(1-2^{-18})N^2>N^2>\mathcal S(N)\). Combining this with the exact scan proves that \(\mathcal F(N)>\mathcal S(N)\) for every integer \(N\ge 8192\).

\begin{table}[htbp]\centering\small
\caption{Selected hyperbolic-cycle lower bounds at exactly \(N\) terms and degree at most \(N-1\). The folded field has \(r\) stages and actual degree \(D_{\rm f}\). The last two columns use the stationary turn and separate strips with \(k=\lfloor N/2\rfloor\).}\label{tab:fold}
\setlength{\tabcolsep}{5pt}
\begin{tabular}{@{}rrrrrr@{}}\toprule
\(N\)&\(r\)&\(D_{\rm f}\)&\(\mathcal F(N)\)&\(\mathcal S(N)\)&\(\mathcal T(N)\)\\\midrule
\inputrows{tex/fold-rows.tex}
\bottomrule\end{tabular}
\end{table}

At \(N=1024\) the separate-strip choice \((i,b,k)=(6,6,512)\) gives \(c_i=1089\). At \(N=8192\), \((8,14,4096)\) gives \(c_i=14{,}081\). Each uses at most \(N-1\) positions before filling and reaches degree \(N-1\). Thus \(\mathcal T=1{,}115{,}136\) and \(115{,}351{,}552\) respectively. At \(8192\) the folded field has already passed the stationary turn, while the separate strips still give more cycles. Allowing \(k\) to vary gives a larger strip count. The triple \((8,15,3841)\) yields \(116{,}036{,}610\) at the same budget (supplement \S\othernumber{sup:lienard-product}{S1.3}~\cite{LuElectronic}).

Among the fifteen budgets \(N=2^j\), \(6\le j\le20\), the first at which \(\mathcal F(N)>\mathcal T(N)\) is \(2^{20}\), but at \(N=2^{21}-1\) one has \(\mathcal F(N)<\mathcal T(N)\). A comparison on the same grid conditional on the seed theorem of an unrefereed Li\'enard preprint appears in Section~\paperlink{low-N-supplement}{subsection.1.4}{sec:seed-replication}{S1.4} of the supplement \cite{LuElectronic}.

For \(48\le N<2^{16}\) the prescribed allocation uses \(q_j=2,\ \alpha_k=4\) throughout. In this case \(J-c<0\) for every \(L\ge1\), by writing \(L\) modulo four in \eqref{eq:Bqalpha}. Only minima and incoming ends contribute. The first stage with \(\alpha_k=8\) occurs at \(r=9\), or \(N=2^{16}\). At stage zero there \(L=16{,}376\), \(J=12{,}279\) and \(c=4{,}093\), giving \(8{,}186\) guaranteed crossings on the outgoing end of each two-well cluster. The first four-well cluster occurs at \(r=22\), or \(N=2^{29}\). The sums in Section~\ref{sec:sum} include growing clusters.

\begin{proof}[Verification of the diagonal allocation]
At \(r=1\), \(q_{\max}=2\), \(N-18R\ge12\). For \(r\ge2\), \(N\ge128R\) and \(q_{\max}\le R\). Hence \(X_*\ge2q_{\max}\), \(s\) is even, \(X_*\le\lfloor N/2\rfloor\), and \(T\ge18R\). Since \(12R\le\lfloor N/2\rfloor\), \(p\ge12R,\ d\ge6R,\ p\le\lfloor N/2\rfloor\) and \(d\le\lceil T/2\rceil\le\lfloor(N-1)/2\rfloor\). The \(d\) slots include exactly the remainder left after choosing \(p,X_*\). Theorem~\ref{thm:finite-folded}, with \(a=2\), produces the field with \(\mathcal F(N)\) cycles. Its general hypotheses require \(p\ge6aR\ge24\) even for a single fold, so \(D\ge47\) and \(N\ge48\) is necessary for that theorem under the linear cap.
\end{proof}
